\documentclass[11pt, a4paper]{article}

\usepackage[T1]{fontenc}
\usepackage[utf8]{inputenc}
\usepackage{tgpagella}
\usepackage{microtype}
\usepackage[spanish, es-noshorthands]{babel}
\addto\captionsspanish{\renewcommand{\tablename}{Tabla}}

\usepackage{amsmath, amssymb, amsfonts}
\usepackage{graphicx}
\usepackage{booktabs}
\usepackage[most]{tcolorbox}
\usepackage[margin=2.5cm]{geometry}
\usepackage{fancyhdr}

\pagestyle{fancy}
\fancyhf{}
\setlength{\headheight}{16pt}
\lhead{\textbf{UCB Sede Tarija} | Ing. Mecatrónica}
\chead{Robótica (IMT-342)}
\rhead{Semana 5 - Sesión 2}
\lfoot{Prof: M.Sc. Ing. Bernardo Quiroga Turdera}
\rfoot{Hoja de Ejercicios}
\renewcommand{\headrulewidth}{0.4pt}
\renewcommand{\footrulewidth}{0.4pt}

\definecolor{ucbblue}{RGB}{0, 51, 102}

\begin{document}

\begin{center}
    \Large\textbf{\textcolor{ucbblue}{UNIVERSIDAD CATÓLICA BOLIVIANA "SAN PABLO"}}\\
    \vspace{0.2cm}
    \normalsize \textbf{Hoja de Ejercicios: Consolidación de Cinemática Directa} \\
    (Resolución Independiente y Diagnóstico)
\end{center}

\vspace{0.4cm}

\section*{Sección A: Problema Diagnóstico 1[cite: 5]}
Un manipulador planar 3R posee la siguiente tabla geométrica (Convención D-H estándar):

\begin{center}
    \begin{tabular}{c c c c c}
    \toprule
    $i$ & $\theta_i$ & $d_i$ [m] & $a_i$ [m] & $\alpha_i$ \\ \midrule
    1 & $q_1$ & 0 & 0.30 & 0 \\
    2 & $q_2$ & 0 & 0.20 & 0 \\
    3 & $q_3$ & 0 & 0.10 & 0 \\ \bottomrule
    \end{tabular}
\end{center}

Determine la pose del efector final respecto al marco base si las variables articulares se encuentran en la configuración:
\begin{equation*}
    q_1=30^\circ, \qquad q_2=-60^\circ, \qquad q_3=90^\circ
\end{equation*}

\vspace{5cm}

\section*{Sección B: Preguntas de Micro-Refuerzo Correctivo[cite: 5]}
\textit{Responda estas preguntas solo después de la discusión diagnóstica con el docente.}
\begin{enumerate}
    \item Para la fila paramétrica $(\theta, d, a, \alpha) = (q, 0.20, 0, -\pi/2)$, identifique el desplazamiento en el eje local $Z$ y el ángulo de torsión (twist) del eslabón.
    \item Dado el conjunto de matrices de transformación ${}^0T_1, {}^1T_2, {}^2T_3, {}^3T_4$, construya la expresión correcta para evaluar ${}^0T_4$.
    \item Dada la siguiente matriz ${}^0T_2$, extraiga analíticamente el vector de posición ${}^0p_2$:
    \begin{equation*}
        {}^0T_2 = \begin{bmatrix} 0&0&1&0.20 \\ 1&0&0&-0.10 \\ 0&1&0&0.30 \\ 0&0&0&1 \end{bmatrix}
    \end{equation*}
\end{enumerate}

\newpage
\section*{Sección C: Problema Representativo 2 (Espacial)[cite: 5]}
Considere un manipulador serial 3R con la siguiente tabla D-H estándar:

\begin{center}
    \begin{tabular}{c c c c c}
    \toprule
    $i$ & $\theta_i$ & $d_i$ [m] & $a_i$ [m] & $\alpha_i$ \\ \midrule
    1 & $q_1$ & 0.20 & 0 & $+\pi/2$ \\
    2 & $q_2$ & 0 & 0.30 & 0 \\
    3 & $q_3$ & 0 & 0.20 & 0 \\ \bottomrule
    \end{tabular}
\end{center}

Para la configuración $q_1=0^\circ, q_2=90^\circ, q_3=-90^\circ$, complete la siguiente secuencia analítica:

\vspace{0.4cm}
\noindent \textbf{1. Transformaciones Individuales:}
\begin{align*}
    {}^0T_1 &= \underline{\hspace{7cm}} \\[1cm]
    {}^1T_2 &= \underline{\hspace{7cm}} \\[1cm]
    {}^2T_3 &= \underline{\hspace{7cm}}
\end{align*}

\noindent \textbf{2. Cadena de Composición:}
\begin{equation*}
    {}^0T_3 = \underline{\hspace{10cm}}
\end{equation*}

\vspace{0.4cm}
\noindent \textbf{3. Matriz Resultante ($T_{\text{final}}$):}
\begin{equation*}
    {}^0T_3 = \underline{\hspace{10cm}}
\end{equation*}

\vspace{1.5cm}
\noindent \textbf{4. Extracción de Pose:} \\
Posición $p = $ \underline{\hspace{5cm}} \quad Orientación $R = $ \underline{\hspace{5cm}}

\vspace{1.5cm}
\noindent \textbf{5. Verificación (Demuestre por qué es físicamente plausible):} \\
\underline{\hspace{15cm}} \\
\underline{\hspace{15cm}}

\end{document}
